An entire function $f$ is of finite order if there is a positive constant $a$ and an $r_0 > 0$ such that $|f(z)| < \exp (|z|^a)$ for $|z| > r_0$. If $f$ is not of finite order then $f$ is of infinite order.

If $f$ is of finite order then the number $\lambda = \inf \{a : |f(z)| < \exp (|z|^a)\}$ for $|z|$ sufficiently large is called the order of $f$.

Notice that if $|f(z)| < \exp (|z|^a)$ for $|z| > r_a > 1$ and $b > a$ then $|f(z)| < \exp (|z|^b)$. The next proposition is an immediate consequence of this observation.